% =====================================================================
%  Capitolo 1 — Introduzione al corso e al mestiere del game designer
%  Fonte: trascrizione della lezione GD1 (Game Design, primo anno)
%  Pacchetti richiesti nel preambolo del libro:
%    \usepackage[italian]{babel}, \usepackage[T1]{fontenc},
%    \usepackage[utf8]{inputenc} (solo pdfLaTeX), \usepackage{booktabs}
% =====================================================================

\chapter[Il mestiere del game designer]{Il mestiere del game designer: introduzione al corso}
\label{cap:introduzione-game-design}

Questo primo capitolo fa da soglia. Prima di affrontare metodi, framework e
documenti, è necessario capire \emph{in che modo} si impara il game design,
\emph{con chi} lo si fa, \emph{con quali strumenti} e, soprattutto,
\emph{che cosa fa concretamente} un game designer nel suo lavoro quotidiano.
Il capitolo segue l'impostazione della lezione inaugurale del corso: si parte
dal contesto formativo (il corso come simulazione di lavoro), si passa al
metodo di apprendimento (l'\emph{inquiry-based learning}), si elencano
letture e risorse di riferimento, e si arriva infine a una prima definizione
operativa del ruolo del designer, dei limiti entro cui lavora e delle sue
possibili specializzazioni.

Una premessa terminologica: il game design è una disciplina giovane e quasi
tutta la sua letteratura è scritta in inglese. Per questo motivo molti termini
tecnici verranno lasciati in lingua originale (\emph{pipeline},
\emph{pitching}, \emph{tuning}, \emph{feature team}, \dots): sono i termini che
si useranno nel lavoro, ed è importante imparare a riconoscerli e a usarli
correttamente fin da subito.

% ---------------------------------------------------------------------
\section{Il contesto: un corso come simulazione di lavoro}
\label{sec:simulazione-lavoro}

\subsection{Il principio fondamentale}

Il primo concetto da fissare è che il corso \textbf{si presenta come una
simulazione del lavoro}. Dal momento in cui si entra in aula si viene trattati
come \emph{potenziali professionisti dell'industria}: non soltanto studenti,
ma lavoratori effettivi di uno studio (l'accademia), con una gerarchia, dei
referenti, dei compiti e delle scadenze.

Ne derivano alcune conseguenze pratiche:
\begin{itemize}
  \item il docente non assegna un ``compitino'' da risolvere per accumulare
        punti o presenze: ogni \emph{task} è pensato per far dedicare tempo
        all'apprendimento di qualcosa di concreto;
  \item non consegnare un task non danneggia il docente (che ``vive ancora
        domani e viene pagato ugualmente''), ma lo studente, che perde
        l'occasione di mettere le mani su un argomento e diventare più abile;
  \item docenti e studenti devono trattarsi \textbf{da colleghi}, perché è
        molto probabile che un domani lo diventino davvero.
\end{itemize}

\subsection{Perché l'impegno conta più che mai}

Il lavoro del game designer è descritto come \emph{molto romantico, ma anche
molto duro}, e sempre più duro col passare degli anni. Durante il periodo del
Covid l'industria ha visto grandi investimenti; quei ``rubinetti'' sono stati
chiusi da tempo, il rapporto con i publisher è diventato più complesso e le
opportunità lavorative non sono più quelle di tre o quattro anni fa. Questo
significa:
\begin{itemize}
  \item più competizione per ogni posizione;
  \item aspettative più alte: ci si aspetta che i candidati siano ancora più
        preparati e sicuri delle proprie competenze;
  \item un unico modo per diventarlo: \textbf{lavorare}. Più tempo si dedica
        oggi ad apprendere il mestiere, più facilmente si sarà assumibili
        domani e più agevole sarà il lavoro di tutti i giorni.
\end{itemize}

\subsection{Sbagliare qui per non sbagliare fuori}

Il mantra del corso è: \textbf{quanto più provate qui, anche in modo goffo e
sbagliando, tanto meno sbaglierete là fuori}. In accademia, per quanto la
simulazione sia realistica, nessuno toglie lo stipendio e nessuno licenzia da
un momento all'altro. Nell'industria sì, e un errore può costare molto denaro
all'azienda, ancora di più se l'azienda è la propria.

\subsection{Il valore del fallimento}

Il fallimento, se affrontato con la mentalità corretta, è probabilmente
\emph{la cosa più vantaggiosa} che si possa fare durante la formazione:
\begin{itemize}
  \item un successo (``avete superato questa deadline'') conferma le proprie
        abilità, ma insegna relativamente poco;
  \item dopo $n$ fallimenti si impara moltissimo, e i successi arrivano;
  \item è fondamentale non demoralizzarsi: il feedback negativo non è mai un
        giudizio sulla persona, ma su ciò che è stato consegnato.
\end{itemize}

Il docente riconosce apertamente che dedicare tempo ed energia a un progetto e
vederlo rifiutato, o vederlo non uscire mai, è demotivante: è successo a lui e
lo ha visto accadere anche su progetti molto grandi, con team numerosi. Il
punto però non è ciò che si è ottenuto materialmente, ma \textbf{quanto si è
cresciuti}. Da ogni esperienza ci si portano a casa due cose:
\begin{enumerate}
  \item \textbf{l'esperienza}: quanto si è diventati bravi, veloci,
        praticamente abili nell'applicare un framework, quanto si è cresciuti
        come designer;
  \item \textbf{il rapporto umano}: nel lavoro si ``suda sangue'' fianco a
        fianco con altre persone, e questo crea legami molto forti che
        portano avanti l'industria nel senso più positivo del termine: nascono
        progetti privati, ci si porta da un'azienda all'altra, si formano team
        con basi solide.
\end{enumerate}

Va evitato anche l'eccesso opposto: la cultura del \emph{grinding} si è
appropriata di questi discorsi in modo distorto. Non bisogna diventare robot.

% ---------------------------------------------------------------------
\section{Organizzazione del corso e gerarchia}
\label{sec:gerarchia-corso}

\subsection{La struttura gerarchica}

La struttura dei referenti del corso ricalca volutamente una
\textbf{gerarchia tipica dell'industria}:
\begin{itemize}
  \item al vertice c'è il \textbf{direttore del corso} (Diego Ricciuti), che
        corrisponde al \emph{creative director}: supervisiona la visione del
        progetto (il concetto di ``visione'' verrà approfondito più avanti);
  \item sotto di lui ci sono i \textbf{referenti di dipartimento}, uno per
        disciplina, con un corrispettivo per ciascuna sede (Milano e Roma);
  \item in fondo ci sono gli studenti, che occupano il primo gradino di
        un'azienda: l'\emph{intern} (stagista) o il \emph{junior}, il passo
        che precede il lavoratore ``regolare''.
\end{itemize}

\begin{table}[htbp]
  \centering
  \begin{tabular}{@{}p{0.34\textwidth}p{0.29\textwidth}p{0.29\textwidth}@{}}
    \toprule
    \textbf{Disciplina} & \textbf{Milano} & \textbf{Roma} \\
    \midrule
    Direzione (\emph{creative director}) & \multicolumn{2}{p{0.58\textwidth}}{Diego Ricciuti} \\
    Game design      & Mattia Podini        & Giacomo ``Jack'' Danna \\
    Level design     & Stefano Zocchi       & Andrea Dispoli \\
    Narrative design & \multicolumn{2}{p{0.58\textwidth}}{Daniele Fusetto (lezioni condivise), Giacomo Masi} \\
    Programmazione (Unreal) & Alessandro Pregnolato & Matteo Velardi \\
    \bottomrule
  \end{tabular}
  \caption{Referenti del corso per disciplina e sede (nomi come riportati in lezione).}
  \label{tab:referenti}
\end{table}

Per il primo anno, il referente diretto degli studenti di game design è il
docente di game design: è a lui che ci si rivolge per dubbi e domande. Altre
figure utili sono il direttore didattico (severo, ``giustamente severo'', da
contattare solo per questioni particolarmente importanti), l'amministrazione
(contratti e documentazione) e la segreteria.

\subsection{I canali di comunicazione}

\begin{description}
  \item[Discord] è il canale per le comunicazioni ufficiali del corso e per
        le domande che non richiedono urgenza immediata; può avere qualche
        ritardo.
  \item[E-mail] è il canale per tutto ciò che è formale o ``ancora più
        ufficiale'' e per le comunicazioni rapide e immediate.
  \item[Telegram] esistono residui di vecchi canali precedenti a Discord, da
        evitare.
\end{description}

Le domande che valgono per tutti vanno fatte in classe (anche da un portavoce
a nome del gruppo); quelle singole e personali via Discord o, se formali, via
e-mail. Se il docente non risponde a una domanda c'è un motivo, ad esempio
ragioni legali, e lo dirà esplicitamente.

\subsection{I tre moduli e la settimana tipo}

Le tre discipline principali del primo anno (e non solo) sono:
\begin{itemize}
  \item \textbf{Game Design} (tendenzialmente il lunedì);
  \item \textbf{Level Design} (tendenzialmente il mercoledì);
  \item \textbf{Programmazione} (tendenzialmente il venerdì).
\end{itemize}
Nel primo mese le lezioni sono tutte di game design (lunedì, mercoledì e
venerdì). Sono previste anche lezioni di narrativa, ma il piano formativo si
aggiorna ogni anno: le ore di narrativa sono state ridotte a favore di quelle
di programmazione. Le poche lezioni di narrativa daranno quindi dei
\emph{cenni}, non la preparazione di un narrative designer. Ciò che serve
davvero è imparare \textbf{come comunicano gli altri dipartimenti} e come
lavorare con persone che fanno un lavoro diverso ma complementare al proprio.

\subsection{Le due anime delle lezioni di game design}

\begin{description}
  \item[Design methods (metodologie)] lezioni prevalentemente teoriche, ma con
        una buona componente pratica, in cui si studiano metodologie di design
        e i cosiddetti \emph{framework} dell'industria.
  \item[Documentazione e pitching] il cuore del lavoro quotidiano del
        designer.
\end{description}

\paragraph{La documentazione.}
Il docente si definisce ``il più grande nerd della documentazione'' che si
possa incontrare. È una parte dura, perché a pochi piace documentare, ma il
lavoro del designer \emph{è principalmente questo}: scrivere documenti e, da
quei documenti, passare alla pratica e all'implementazione. Chi non sa
documentare bene parte male. Quando non si è presenti, il documento è il
proprio \textbf{biglietto da visita}: qualcun altro lo leggerà, quindi va
scritto nel modo più chiaro e comprensibile possibile.

Un esempio concreto del perché la documentazione conta: CD Projekt Red ha
sviluppato un proprio engine proprietario, ma buona parte della sua
documentazione è stata ``buttata all'acqua''. Chi deve usarlo o lo sa già
usare, o ha competenze di un altro engine da trasferire, oppure è in
difficoltà. Si documenta \textbf{perché altri possano leggere}.

\paragraph{Il pitching.}
Il \emph{pitching} è l'esposizione, la presentazione orale. La parte creativa
del team presenta spesso:
\begin{itemize}
  \item le proprie idee ai colleghi;
  \item il gioco a convention ed eventi;
  \item il pitch del gioco agli investitori, per ottenere i fondi necessari a
        realizzarlo.
\end{itemize}
Saper parlare e comunicare chiaramente le idee è fondamentale. In sintesi:
\emph{il game designer parla}, e bisogna farsene una ragione.

\subsection{Altre componenti del corso}

\begin{description}
  \item[Gioco critico e playtest] la classe viene divisa in due gruppi: uno
        gioca, l'altro osserva. Chi osserva non applaude ``come le
        scimmiette'': annota, segna e analizza il \emph{playtest}, esattamente
        come si fa nell'industria. Il playtest fa parte delle mansioni del
        designer, ma esiste anche come specializzazione autonoma.
  \item[Pipeline] termine che indica tutto ciò che riguarda i processi
        formalizzati dell'industria; di solito si intende la \emph{pipeline di
        produzione}: come si divide il team, qual è la sua gerarchia, come il
        \emph{producer} gestisce le deadline. In breve: \emph{come} e
        \emph{quando} si fanno le cose. Al primo anno è meno rilevante (si
        dovrà ``nuotare'' un po' da soli); dal secondo anno si vedranno le
        metodologie effettivamente applicate ai progetti digitali. Si
        scoprirà quanto sia difficile lavorare a un progetto digitale, e
        ancor più fisico, senza un piano di sviluppo comprovato, come quelli
        usati per qualsiasi software in qualsiasi parte del mondo.
  \item[Engine] il motore di riferimento è \textbf{Unreal Engine}
        (sezione~\ref{sec:engine}).
  \item[Programmazione] principalmente tramite \emph{Blueprint} in Unreal,
        più eventuali sistemi di versionamento; il corso parte da zero.
  \item[Progetti pratici] probabilmente la cosa più importante. Al primo
        anno i progetti di gruppo restano circoscritti al corso di design,
        perché mancano ancora alcune nozioni (come comunicare, affrontare
        problemi, dare e ricevere feedback). In seguito si allargheranno a
        progetti digitali, molto più difficili di quelli fisici, con membri
        di altri team: un rapporto non unilaterale, ma ``covalente''.
\end{description}

% ---------------------------------------------------------------------
\section{Il metodo: l'inquiry-based learning}
\label{sec:inquiry}

\subsection{Che cos'è}

Il corso è costruito secondo l'\textbf{inquiry-based learning}
(apprendimento basato sull'indagine), un approccio pedagogico nato intorno
agli anni Sessanta, che si richiama a pensatori come Dewey e Piaget. In
generale, le discipline da cui il design attinge sono quasi sempre la
sociologia, la psicologia e, talvolta, l'economia.

Il funzionamento è un ciclo:
\begin{enumerate}
  \item il docente fornisce metodologie e framework dell'industria durante
        la lezione;
  \item lo studente, fuori dalle ore di lezione, lavora sui \emph{task}:
        mette a terra le nozioni, le maneggia, le assorbe meglio;
  \item da questo lavoro nascono nuovi dubbi;
  \item i dubbi tornano in aula: lo studente chiede, il docente risponde
        (ammesso che dare la risposta sia utile, e ammesso che una risposta
        esista).
\end{enumerate}

\subsection{Non sempre esiste una risposta giusta}

Non tutto ciò che si studia ha ``la risposta migliore possibile in assoluto
nel vuoto''. Esistono metodologie, limiti e framework che aiutano il lavoro,
ma non sono fissi: cambiano da azienda ad azienda, da lavoro a lavoro, da
progetto a progetto, e a volte vengono stravolti. Sta alla
\textbf{sensibilità del designer} capire quale strumento o framework usare in
un dato momento, e questa sensibilità si impara solo con l'esperienza. Nel
lavoro nessuno dirà ``usa questa cosa'' o ``è meglio quest'altra'': sarà il
designer a scegliere.

Anche l'esperienza del docente ha confini: sei-sette anni di lavoro, un modo
personale di approcciare il design e un certo insieme di progetti (ad esempio
nessun MMO). Oltre quei confini si cercano risposte nella ricerca, studiando
come aziende con esperienza specifica (ad esempio Riot) hanno affrontato
problemi simili e hanno condiviso le loro soluzioni, sempre valutando quali
metodi siano adatti alla dimensione del team, al tipo di progetto e al tempo
disponibile.

\subsection{Sensibilità, non memoria}

Non interessa imparare a memoria un framework o la teoria di un autore. Certe
terminologie vanno conosciute, ma per ragioni di comunicazione e chiarezza.
Ciò che conta è sviluppare la \textbf{capacità di farsi domande}, di farle
agli altri, di fare ricerca e di trovare una propria risposta. È esattamente
ciò che si fa nel lavoro: davanti a qualcosa che non si sa fare, l'abilità
del designer sta nel cercare come qualcun altro ha risolto un problema simile.

\subsection{Le lezioni come corsi di aggiornamento}

Le lezioni frontali simulano i \emph{corsi di aggiornamento} aziendali: le
sessioni che un lead di design tiene al suo team di design (o il lead di
level design al suo team, e così via). Alla fine dell'aggiornamento c'è quasi
sempre un lavoro da fare, un obiettivo da portare a terra:
\textbf{l'obiettivo della settimana}. Si lavora di settimana in settimana, e
il perché verrà chiarito in seguito.

Struttura tipica di una lezione:
\begin{enumerate}
  \item breve sessione di domande e risposte sulla lezione precedente;
  \item analisi di eventuali esercizi particolarmente critici;
  \item lezione vera e propria;
  \item discussione di gruppo, se c'è tempo;
  \item assegnazione del task della settimana.
\end{enumerate}
Il docente tende a ``parlare e straparlare'', ad aprire parentesi e a saltare
da un argomento all'altro: l'unico modo per fermarlo è che siano gli studenti
a interromperlo con domande.

\subsection{La gerarchia della competenza}

L'inquiry-based learning può essere rappresentato come una scala di
competenza:
\begin{enumerate}
  \item la sola lezione frontale porta a \textbf{conoscenze incomplete},
        perché manca la parte pratica e di ricerca;
  \item il lavoro a casa (task, ricerca) porta a una competenza \textbf{più
        completa}, non sbagliata in sé, ma ancora parziale e complementare ad
        altre;
  \item la \textbf{condivisione} dei task e dei feedback con gli altri porta
        idealmente alla competenza più completa possibile;
  \item il picco della competenza è la sua \textbf{applicazione nei
        progetti} (fisici o digitali, di gruppo): ``con le cose che avete
        imparato e ricercato, riuscite a fare un gioco?''. Quando la risposta
        diventa positiva, o comunque si inizia a navigare nei processi di
        design per realizzarne uno, quella è la vittoria.
\end{enumerate}
Non si parte subito creando giochi: si parte dall'analisi, dal design e dalle
regole generali di design, prima ``nel vuoto'', poi contestualizzando i
framework, infine unendo i puntini nei progetti.

Rimarranno comunque delle lacune, per tutti, docente compreso, perché il game
design \textbf{non è una scienza esatta}, anzi, a rigore non è affatto una
scienza. Prende in prestito metodologie dalle discipline scientifiche, ma
contiene una componente imprevedibile e incalcolabile: \textbf{la
produzione}. Si possono applicare perfettamente le regole di un framework, e
quel framework può comunque non essere adatto al gioco, al proprio modo di
sviluppare o al tempo di produzione disponibile (alcuni framework richiedono
una ricerca e un'attenzione che il tempo non concede).

C'è poi la componente creativa, la parte ``romantica'' del lavoro: esiste, ma
è \textbf{metà del lavoro}. L'altra metà è ricerca, metodo scientifico al
100\%. Dagli stessi dati $n$ designer diversi, con sensibilità diverse,
possono ottenere $n$ risposte diverse. Non tutte saranno valide: vanno
applicate, messe a terra, e questa scrematura fa parte del lavoro.

\subsection{I livelli dell'apprendimento}

Un secondo modello, a gradini, descrive i livelli di apprendimento (è la
cosiddetta tassonomia SOLO):
\begin{description}
  \item[Pre-strutturale] il docente dà il compito e lo studente lo risolve,
        possibilmente il prima possibile. È esattamente ciò che \emph{non} si
        vuole.
  \item[Uni-strutturale] si affronta una sola disciplina: è il primo passo
        necessario, ed è il motivo per cui nel primo mese ci sono solo lezioni
        di game design.
  \item[Multi-strutturale] entrano in gioco altri dipartimenti: level design,
        game design, programmazione e così via.
  \item[Relazionale] si capisce che alcune cose valgono per più
        dipartimenti: alcuni framework di narrativa derivano dal design o
        possono essere usati per fare design; le regole di documentazione si
        applicano anche ad altri dipartimenti; le regole di level design si
        mescolano con quelle di game design. Dopotutto, per fare un gioco
        servono molti dipartimenti.
  \item[Astratto esteso] il punto finale: muoversi autonomamente e capire
        cosa usare e quando, sulla base della propria esperienza. L'unico modo
        per arrivarci è fare pratica e fare progetti.
\end{description}

% ---------------------------------------------------------------------
\section{Task, feedback e regole del corso}
\label{sec:task-feedback}

\subsection{I task come opportunità}

I task non sono strettamente obbligatori (esiste un numero minimo di
consegne), ma sono \textbf{opportunità}. Quanto e come ci si lavora dipende
dallo studente. Un task fatto in cinque minuti, o in cui si è stati ``furbi'',
non verrà letto né riceverà feedback.

Domande da porsi davanti a ogni task:
\begin{enumerate}
  \item \textbf{Perché} mi è stato assegnato? Qual è l'obiettivo?
  \item Come ci lavoro per ottenere un risultato \textbf{che abbia senso}? Se
        la ricerca e i dati indicano una direzione, è inutile andare in
        un'altra.
\end{enumerate}

\subsection{Quanto lavorare: un calcolo indicativo}

Se il corso è una simulazione di lavoro, vale il carico di un lavoro normale:
\begin{itemize}
  \item un lavoro a tempo pieno comprende da 30 a 45 ore settimanali; si
        considera una media di \textbf{40 ore};
  \item si sottraggono circa 9 ore di lezione: restano $\approx 31$, diciamo
        \textbf{30 ore};
  \item con circa 3 task a settimana: $30 / 3 = $ \textbf{circa 10 ore per
        task}.
\end{itemize}
Il numero può variare, ma rende l'idea: dedicare un'ora a un task e poi fare
altro non fa guadagnare nulla a nessuno. Lo studente ``butta soldi'', il
docente risparmia tempo leggendo il task in un minuto (o non leggendolo).
Dedicando invece attenzione, ricerca, cura nella formattazione e
sperimentazione, e facendolo con cognizione di causa, si impara.

\subsection{Niente voti, solo feedback}

Ad eccezione del corso di programmazione, non ci sono voti (né ``ottimo'', né
10, né 100). Ci sono \textbf{feedback}:
\begin{itemize}
  \item in game design, generalmente sulla documentazione;
  \item in level design, sui progetti pratici.
\end{itemize}
Il docente indica cosa va bene e cosa va modificato, concentrandosi in
particolare su ciò che va cambiato. Il feedback \textbf{non è mai un giudizio
sulla persona}, ma su ciò che è stato consegnato nel concreto.

\subsection{La condivisione dei task}

Tutti i task consegnati vengono \textbf{condivisi} con la classe, e possono
essere letti solo dopo aver ricevuto i feedback del docente. Si cresce così
non solo con i propri feedback, ma anche con quelli dati agli altri. Regola
importante: ciò che il docente dice a uno studente \textbf{vale per tutti}. Se
a una persona viene detto ``questa cosa non va mai più fatta'', da quel
momento vale per l'intera classe.

\subsection{Non copiare}

Copiare i task degli altri è una delle cose peggiori che si possano fare, non
per una questione morale, ma perché annulla il senso del percorso. Si può
dedicare a ogni task il tempo e l'attenzione che si ritengono opportuni, ed è
normale che col tempo ci si specializzi (preferendo, ad esempio, il level
design o la programmazione tecnica) e si investa meno dove si sente meno
necessità. In questi casi l'alternativa corretta è:
\begin{itemize}
  \item non consegnare;
  \item studiare ciò che hanno fatto gli altri;
  \item eventualmente riconsegnare in seguito.
\end{itemize}
Ci saranno momenti in cui si verrà sommersi dalle cose da fare: bisogna
imparare a scegliere.

\subsection{Il ``non leggo''}

Il docente ha piena facoltà di non dare alcun feedback a un task (il
cosiddetto \textbf{``non leggo''}) se:
\begin{itemize}
  \item non si è lavorato con attenzione;
  \item non si rispettano le indicazioni date;
  \item è stata violata una regola estremamente importante.
\end{itemize}
Un task ``non letto'' equivale a un task non consegnato e non entra nel
conteggio. Attenzione: un task \emph{perfetto} e un task che \emph{fallisce}
nel comprendere o applicare un framework sono \textbf{ugualmente validi},
purché ci si sia impegnati e si sia dedicato tempo, pazienza e ``sudore''.

\paragraph{Riconsegna.} Un task non letto, in ritardo o che si vuole
migliorare può essere riconsegnato e verrà conteggiato, a patto che:
\begin{itemize}
  \item i problemi segnalati siano stati risolti;
  \item sia stato fatto ulteriore lavoro, integrando anche tutto ciò che nel
        frattempo è stato detto agli altri colleghi nei loro feedback.
\end{itemize}
Il task va sviluppato con le competenze che si hanno nel momento della
riconsegna, non con quelle del momento della prima consegna.

\subsection{Requisiti per superare l'anno}

\begin{itemize}
  \item consegnare almeno l'\textbf{80\% dei task} assegnati (soglia
        malleabile);
  \item partecipare almeno all'\textbf{80\% delle lezioni}, in presenza o
        da remoto, preferibilmente in presenza; i ritardi vengono valutati
        caso per caso;
  \item comportarsi da persone corrette e rispettose (è stata firmata una
        carta dello studente): non urlare, non alzare le mani, non
        vandalizzare, non compiere atti illegali.
\end{itemize}
Impegnarsi di più porta inoltre a essere selezionati per varie opportunità.

\subsection{Scadenze}

\begin{itemize}
  \item la scadenza di un task è \textbf{sempre prima dell'inizio della
        lezione} successiva indicata;
  \item se la classe, come team, mantiene comportamenti scorretti, il docente
        può \textbf{anticipare le deadline} per tutti.
\end{itemize}

\subsection{Accessibilità e DSA}

Gli studenti con disturbi specifici dell'apprendimento o neurodivergenze
(dislessia, ADHD, autismo, ecc.) sono invitati a segnalarlo. Quando si
affronteranno le regole di formattazione e documentazione (errori
grammaticali, mancanze, incoerenze) se ne terrà conto e si potranno trovare
dei compromessi, pur ricordando che nel mondo del lavoro nessuno farà
``sconti''. La documentazione può essere inviata all'amministrazione o alla
segreteria, che la comunicherà a tutti i docenti, ma il consiglio è parlarne
anche direttamente con ogni singolo docente.

% ---------------------------------------------------------------------
\section{Le competenze trasversali}
\label{sec:soft-skill}

\subsection{I giochi si fanno in gruppo}

Si possono fare giochi da soli, ed è un'impresa ammirevole; ma lavorando da
soli \textbf{tutto ciò che si guadagna in libertà si perde in produttività}.

\paragraph{Il caso Stardew Valley.} \emph{Stardew Valley} è stato realizzato
da una sola persona, Eric Barone, in anni di sviluppo e di sacrifici: faceva
un altro lavoro per mantenersi (in un cinema) e di notte programmava, faceva
l'arte e componeva la musica. Il risultato è stato straordinario e lo ha reso
ricco, ma si tratta di un caso estremamente raro. Quanti giochi simili, nati
nelle stesse condizioni, sono falliti? Molti di più.

\paragraph{Il rischio d'impresa.} Lavorare da soli non è sbagliato di per sé
e ha aspetti molto positivi: si ha il privilegio di creare esattamente ciò che
si vuole. Ma ci si fa carico anche di tutto il \textbf{rischio d'impresa}: se
si lavora a un gioco per cinque anni e poi il gioco va male, resta il
vantaggio di averlo pubblicato, ma il dopo è sperare di vivere di quella
release o di finanziare con essa il gioco successivo. Le grandi aziende hanno
gli stessi problemi, ma li distribuiscono su più persone e più responsabilità.

L'obiettivo di un professionista è \textbf{saper lavorare in diverse
condizioni}, e la più comune è quella del dipendente (o del freelance /
libero professionista con partita IVA): si viene assunti per le proprie
capacità da un'azienda, piccola o grande che sia, e si lavora con altre
persone.

\subsection{Essere persone con cui si vuole lavorare}

Se le persone non vogliono lavorare con voi, quello è \textbf{il primo motivo
di licenziamento}. E se non si viene licenziati, il clima nell'azienda
diventa terribile: la produttività crolla, la motivazione crolla. Le
\emph{soft skill}, cioè l'insieme di abilità che rendono una persona
piacevole con cui lavorare, vanno coltivate e diventano sempre più importanti
col passare degli anni (relativamente meno ora, molto di più quando i team
saranno più grandi). Anche chi è un ``lupo solitario'' deve imparare a:
\begin{itemize}
  \item lavorare con altre persone;
  \item comunicare con altre persone;
  \item trovare un punto in comune da cui partire per fare le cose.
\end{itemize}

\subsection{Indipendenza e proattività}

Sono le altre due parole chiave del corso:
\begin{description}
  \item[Indipendenza] spesso si verrà lasciati da soli, non per cattiveria ma
        per sfida: quando si cerca una risposta da soli, ciò che conta non è
        tanto la risposta quanto \textbf{il processo}. Una risposta data dal
        docente non vale niente.
  \item[Proattività] chiedere, domandare, agire da soli senza aspettare
        direttive, perché spesso quelle direttive non ci sono, a maggior
        ragione nei team piccoli.
\end{description}
Tutto questo si concretizza nel parlare, nel presentare e nel fare
\emph{post-mortem} (l'analisi a posteriori di un progetto concluso, che verrà
approfondita più avanti), e in generale nell'attenzione alla crescita
personale.

\subsection{Fare domande}

Il corso insiste sulle domande: alzare la mano, intervenire anche ``a
bomba'', fermare il docente ogni volta che qualcosa non è chiaro. Se una
persona non conosce qualcosa, molto probabilmente c'è almeno un'altra persona
che non la conosce. Il docente può dare per scontata la conoscenza di termini
che non lo sono, ed è del tutto normale, soprattutto all'inizio.

% ---------------------------------------------------------------------
\section{Strumenti: engine, Unreal e materiali}
\label{sec:engine}

\subsection{Unreal Engine e gli altri motori}

Il motore di riferimento del corso, usato in level design, è \textbf{Unreal
Engine}: ``un gran mattone'', ma che dà a un designer molte possibilità. Il
vantaggio principale è che permette a un designer non pratico di
programmazione di \textbf{prototipare tramite i Blueprint} (il sistema di
scripting visuale di Unreal). Questo non impedisce di passare a Unity o ad
altri engine, e viceversa: il docente stesso ha studiato e lavorato a lungo
con Unreal, ma la sua ultima esperienza professionale (come level designer) è
stata su Unity, un engine che non conosceva, a dimostrazione che il salto da
un engine all'altro è possibile.

Panorama dei motori citati:
\begin{itemize}
  \item \textbf{Unreal Engine} e \textbf{Unity}: i due più diffusi, con
        modelli di sorgente ``più o meno'' aperto;
  \item \textbf{Godot}: engine open source;
  \item \textbf{CryEngine}: proprietario di Crytek;
  \item \textbf{Snowdrop}: engine di Ubisoft;
  \item Rockstar usa un proprio engine proprietario;
  \item alcune aziende usano versioni \textbf{personalizzate} di engine
        esistenti: ad esempio Milestone, nei giochi di corse, usa un Unreal
        customizzato;
  \item CD Projekt Red ha un engine proprietario (con i problemi di
        documentazione già citati).
\end{itemize}
In sintesi: ogni azienda può avere un engine proprietario oppure una versione
modificata di un engine commerciale, ma si parte sempre da una base.

\subsection{Programmazione: cosa ci si aspetta da un designer}

Non è richiesto alcun bagaglio pregresso di programmazione: il corso parte da
zero. Gli studenti sono stati selezionati (il corso è a numero chiuso) per
competenze analitiche e critiche sopra la media. Anche nel mercato del lavoro,
per una posizione di game designer generalmente \textbf{non ci si aspetta
competenza di scripting} (C++, Java, C\#, ecc.), a meno di voler fare il
\emph{technical designer}, e anche lì varia da azienda ad azienda.

\subsection{Versione e preparazione}

Indicazioni pratiche emerse in aula:
\begin{itemize}
  \item scaricare Unreal Engine; la versione consigliata per il corso è
        stata indicata nella \textbf{5.6} (con conferma successiva). La
        stabilità dipende da cosa si deve fare, da quali tool e plugin si
        usano: i problemi nascono spesso dai plugin. Nel passaggio da Unreal 4
        a Unreal 5, ad esempio, il sistema \emph{Landscape} era inizialmente
        inutilizzabile, mentre in Unreal 4 funzionava molto bene;
  \item è consigliato installare tramite \textbf{Epic Games Launcher}, più
        comodo: permette anche di riscattare giochi gratuiti periodici e di
        accedere ai pacchetti di \textbf{Fab} (il marketplace di asset
        acquisito da Epic), che ogni mese offre alcuni asset gratuiti utili
        per esercitarsi nel level design (ad esempio per fare \emph{speed
        level design});
  \item molti pacchetti di UI o di sistemi di interazione sono utili per
        prototipare al volo, ma spesso inutilmente complicati;
  \item un plugin di \textbf{materiali di prototipazione a griglia} (adattivi,
        con metriche, allineabili alla griglia, con materiali già pronti e
        colori modificabili, dal costo di circa 5 euro) è indicato come
        ``salvavita'' per il level designer; se necessario, i pacchetti
        indispensabili vengono forniti;
  \item portare \textbf{sempre il portatile} alle lezioni di level design
        (da novembre, a partire dall'interfaccia dell'engine); alle lezioni
        di game design basta che lo porti almeno metà della classe, per i
        lavori a coppie. Serve comunque qualcosa per prendere appunti
        (``tavolette d'argilla e cannucce, se volete fare i babilonesi'',
        oppure carta e penna).
\end{itemize}
Chi ha un portatile poco potente dovrebbe comunque provare a installare
l'engine e aprire un template per vedere come gira: all'inizio non si faranno
cose pesanti. Nelle prime lezioni di level design è importante mettere subito
le mani sull'interfaccia: ascoltare dove si trova uno strumento (ad esempio lo
\emph{snapping}) senza provarlo porta a dimenticarlo.

% ---------------------------------------------------------------------
\section{Letture e risorse di riferimento}
\label{sec:risorse}

\subsection{I libri}

\begin{description}
  \item[\emph{The Art of Game Design: A Book of Lenses} (Jesse Schell)] il
        testo \textbf{davvero importante}. Alcune lezioni finali dell'anno si
        baseranno sulle sue terminologie. È utile perché apre la mente sulla
        possibilità di applicare metodologie diverse a giochi diversi, ma
        soprattutto per la \textbf{terminologia} comune. Esiste anche in
        audiolibro (circa 25 ore).
  \item[\emph{Challenges for Game Designers} (Brenda Romero e Ian Schreiber)]
        libro \emph{entry level} con moltissimi esercizi, molto caro al
        direttore del corso.
  \item[\emph{21st Century Game Design} (Chris Bateman)] utile come
        retrospettiva: è un libro visionario che ha previsto metodi e
        soprattutto \emph{archetipizzazioni} dei giocatori tuttora valide nel
        design e nel marketing. Servirà per una delle lezioni più importanti
        dell'anno. È utile leggerlo per vedere quanto aveva previsto, ma è
        utile ``fino a un certo punto''.
  \item[\emph{An Architectural Approach to Level Design} (Christopher
        Totten)] non strettamente necessario, ma di lettura semplice. Ha due
        anime: parte dall'architettura (babilonese, greca, ecc.) e la trasla
        su esempi di level design (come \emph{Doom} e \emph{Team Fortress}).
        Un tempo ne esisteva un riassunto su Gamasutra (oggi
        \emph{Game Developer}).
  \item[\emph{A Theory of Fun for Game Design} (Raph Koster)] utile fino a un
        certo punto: contiene diverse idee superate, che lo stesso Koster ha
        affrontato in vari talk reperibili nel GDC Vault. Verrà ripreso nelle
        prime lezioni, poi messo da parte.
\end{description}

\subsection{Canali e conferenze}

\begin{description}
  \item[GDC (Game Developers Conference)] il suo canale raccoglie talk e
        panel estremamente importanti: è ``uno strumento misterioso che ci
        servirà più tardi''. \textbf{Attenzione}: per molti talk mancano ancora
        le competenze per interpretarli, e alcuni potrebbero fuorviare o
        mettere in testa idee strane. La GDC si concentra spesso su
        produttività e mobile.
  \item[Il \emph{Dragon Speech} di Chris Crawford] il talk del fondatore della
        GDC: appartiene a un passato lontano, ma è molto motivante. Da vedere
        subito.
  \item[Game Maker's Toolkit] canale particolarmente utile per nozioni
        pratiche, con video brevi e alcune interviste.
  \item[Nordic Game] altra conferenza importante, con talk e panel di esperti
        del settore.
\end{description}

\subsection{Associazioni e premi: come leggerli}

Le associazioni di settore (ad esempio IIDEA in Italia, o la IGDA) non sono
il male, anzi: sono un canale importante per affacciarsi all'estero o per
accedere a fondi. Va però ricordato che \textbf{tutelano i propri associati}:
non esiste solo il loro palinsesto, e bisogna valutare quanto spazio occupano
nel panorama della visibilità.

Lo stesso vale per i \textbf{premi} (Game Awards, ``Best Mobile Game'',
``Games for Change'', ecc.): al designer \emph{non interessano}. Che un gioco
abbia vinto ``più di 90 premi'' non dice nulla di utile. Servono a dare
visibilità e lustro ai giochi ed è giusto che esistano, perché ampliano il
panorama e il numero di giocatori. Ma ciò che interessa al designer sono i
\textbf{dati}:
\begin{itemize}
  \item come sono state implementate certe cose;
  \item come è stato fatto il marketing e dove il gioco è stato presentato;
  \item non quante wishlist ha fatto, ma \emph{come} le ha fatte;
  \item in che periodo storico è uscito e quali erano i competitor in quel
        momento.
\end{itemize}

% ---------------------------------------------------------------------
\section{Che cosa fa un game designer}
\label{sec:cosa-fa}

Per arrivare a una prima definizione, la classe ha condotto un brainstorming:
ognuno doveva indicare una cosa che, secondo lui, fa il game designer (senza
ripetere quelle già dette). Ne è emersa la lista seguente, commentata punto
per punto.

\subsection{Progettare meccaniche e struttura del gioco}

Il ruolo più generico è quello del \textbf{core mechanic designer}, detto
anche \emph{gameplay designer} e, un tempo, \emph{black box designer}: si
occupa delle \textbf{meccaniche di gioco}. Si va dal piccolo al grande (in
funzione anche della dimensione dell'azienda):
\begin{itemize}
  \item nel \textbf{piccolo} ci sono le singole meccaniche;
  \item nel \textbf{grande} ci sono i \textbf{sistemi}. Nei giochi
        \emph{sistemici}, con molti sistemi interconnessi (ad esempio
        \emph{GTA}), esistono specializzazioni dedicate a bilanciare e
        supervisionare i sistemi: è il \emph{system design}.
\end{itemize}

\paragraph{Quando un gioco è ``sistemico''?}
Non dipende solo da come le meccaniche interagiscono con il giocatore (le
\emph{dinamiche}, secondo i framework che passano dalle meccaniche alle
dinamiche, che verranno studiati più avanti), ma anche da \emph{quante}
meccaniche ci sono e da quanto interagiscono tra loro. Non diremmo che
\emph{Super Mario} sia un gioco sistemico, eppure contiene dei sistemi: la
linea è sottile. Quando i sistemi aumentano molto, sono complessi e si
influenzano a vicenda, il gioco diventa sistemico. Esempio: in \emph{The
Legend of Zelda: Breath of the Wild} il vento modifica quanto il personaggio
sale o scende; il fuoco può incendiare l'erba; l'erba prende fuoco se colpita
da una freccia infuocata; e così via. Nel mercato moderno i giochi sistemici
sono spesso, ma non sempre, legati a mappe \emph{open world}.

\subsection{Le 3C: camera, character, control}

La base del gioco, da cui generalmente si inizia a lavorare (anche nei
prototipi), è la ``sacra trinità videoludica'' delle \textbf{3C}:
\begin{description}
  \item[Camera] che tipo di inquadratura ha il gioco: prima persona, terza
        persona, isometrica, ecc.
  \item[Character] quanto è grande il personaggio, che tipo di personaggio è
        e, di conseguenza, quali azioni può compiere.
  \item[Control] i controlli: che cosa il personaggio può fare, quali sono i
        \emph{verbi} di gioco.
\end{description}
Le 3C sono definite dal designer, ed esiste anche una specializzazione per
questo lavoro. Esistono varianti: quella del direttore del corso è
\textbf{4C + J}, cioè \emph{camera, character, control, conflict} (conflitto)
e \emph{juice} (il ``succo'', il feedback che rende piacevole l'interazione).
Verranno approfondite più avanti.

\subsection{Il personaggio}

Il designer definisce il personaggio dal punto di vista delle meccaniche
(dimensioni, azioni, capacità), all'interno delle 3C. L'aspetto estetico è
invece competenza di altri dipartimenti.

\subsection{Documentare le informazioni per gli altri dipartimenti}

Assolutamente vero: la documentazione è una delle prime attività del
designer. Più l'azienda è grande, più questo lavoro diventa prominente; più è
piccola, meno è rilevante. In \textbf{pre-produzione}, però, è sempre
fondamentale.

\subsection{Mediatore e leader ``al servizio''}

Il designer è un \textbf{mediatore} nei conflitti interni al team, una sorta
di collante che si relaziona con dipartimenti diversi (non sempre gli stessi,
a seconda della fase di produzione). A volte fa anche da leader, a seconda di
come sono gestite la pipeline e la gerarchia.

Questa leadership è difficile da definire perché \textbf{non è verticale}: è
una leadership \emph{al servizio} (\emph{servant leadership}). Il designer
``è asservito al programmatore'': il suo compito non è arrivare con un'idea e
scaricarla addosso agli altri, né portare la soluzione e lasciare il problema
a qualcun altro, ma \textbf{seguire il processo a basso livello}, sporcandosi
le mani.

\paragraph{I feature team.} Nei team medio-piccoli o medio-grandi possono
esistere sotto-team (\emph{feature team}) dedicati a un obiettivo specifico:
ad esempio un designer incaricato del character o della camera, affiancato da
uno o due programmatori e da un artista. In questi casi il designer
supervisiona il lavoro e gestisce la produzione del sotto-team.

\subsection{Gestire l'avanzamento del prodotto}

Esiste una specializzazione che si occupa solo di questo, il
\textbf{producer}, a un livello gerarchico completamente diverso (e che
potrebbe essere una strada percorribile al termine del percorso di studi).
All'interno di un feature team anche il designer può farlo, ma non è detto.
Può certamente \emph{dare la direzione}: nella sua documentazione può
indicare che per arrivare a uno stato giocabile del prototipo serve un certo
lavoro di arte e di programmazione, per lo stato successivo un altro, e così
via. In ambito produttivo esistono metodologie strutturate come
\textbf{Agile}, che verranno viste in seguito.

\subsection{Presentare le idee}

Il pitching (sezione~\ref{sec:gerarchia-corso}) è richiesto al designer sia
in pre-produzione sia in produzione, in misura che dipende dall'azienda. Ad
esempio Ubisoft lascia ai designer, e in realtà a qualsiasi membro, la
possibilità di \emph{pitchare} un'idea prima di una nuova produzione, perché
possa iniziare a prendere vita.

\subsection{Bozze e prototipi}

Le bozze vanno bene, ma devono essere \textbf{chiare}. C'è un'eterna lotta tra
documentazione e gioco giocabile: se si fa solo documentazione il gioco non
avanza; se si fa solo il gioco manca la documentazione e si perde la visione.
Sono due cose \textbf{complementari}.

Il modo per comunicare un'idea in forma grezza ma giocabile è il
\textbf{prototipo}. Un prototipo può essere qualsiasi cosa: di carta, in
Unreal con Blueprint ``brutti''. \emph{Non deve essere bello, deve essere
giocabile}.

\paragraph{Esempio: il prototipo di carta di Mario + Rabbids.} Durante lo
sviluppo di \emph{Mario + Rabbids}, i programmatori di Ubisoft erano impegnati
su altri progetti (come \emph{Ghost Recon}). I designer realizzarono allora
un \textbf{gioco da tavolo} che riproduceva le meccaniche del gioco, quelle
di uno strategico a turni sul modello di \emph{XCOM}. E funzionava.
L'obiettivo era \textbf{validare un'idea}, validare un gameplay.

Altri modi per rendere chiara un'idea sono i \emph{mock-up} nella
documentazione e le immagini modificate o i \emph{mash-up} di qualsiasi
natura. Lato documentazione l'importante è che sia chiaro; lato prototipo,
che sia giocabile.

\subsection{Fare ricerca per l'idea, e poi smontarla}

Il designer ha una componente creativa importante e fa ricerca per arrivare a
un'idea. Ma ciò che deve poi fare necessariamente è \textbf{smontarla}. Il
suo lavoro è scientifico, esattamente come il \textbf{metodo scientifico}, il
cui obiettivo è \emph{confutare un'ipotesi}: si arriva a un'ipotesi tramite
sperimentazione e poi si cerca $n$ volte di confutarla; se resiste, la si
assume come valida \emph{in un certo contesto} (le regole della fisica
classica, ad esempio, non valgono nella meccanica quantistica).

Allo stesso modo il designer arriva a un'idea (tramite ricerca o chiedendo
agli altri: le idee le hanno tutti, non sono un privilegio del designer),
cerca di smontarla e poi cerca di \textbf{metterla a terra}. Se l'idea non è
applicabile, non è implementabile, finché non si trova un modo per
realizzarla va scartata. Può essere l'idea più bella dell'universo, ma
\textbf{se non può essere messa a terra, non esiste}.

\paragraph{L'originalità.} Serve un'idea originale? Dipende dall'obiettivo e
dalla visione del progetto:
\begin{itemize}
  \item se l'obiettivo è vendere molto, un'idea originale non serve
        necessariamente;
  \item se si vuole creare un nuovo filone, serve un'idea, se non originale,
        almeno interessante: può essere un discorso di marketing o una
        meccanica innovativa, con il rischio di non vendere.
\end{itemize}
Il ruolo del \emph{director} e il concetto di \emph{visione} verranno
approfonditi in seguito.

\subsection{Analizzare i competitor}

Arrivati a un'idea (``vogliamo un certo tipo di combattimento''), si va a
vedere se qualcuno l'ha già fatta e \emph{come}. Le meccaniche non sono
privilegio di nessuno, salvo quando vengono brevettate, come il
\textbf{Nemesis System} di \emph{La Terra di Mezzo: L'Ombra di Mordor}.

\paragraph{Esempio: il combattimento da ridimensionare.} In un progetto
(presso Stormind Games) il team si trovò con un sistema di combattimento
troppo costoso: mancavano le risorse e gli animatori per realizzarlo. Poco
output lato gameplay significava non poter fare qualcosa di complesso e
profondo come in \emph{Elden Ring}; al contrario, il team aveva molti
artisti e poteva puntare su qualcosa di più artistico. Si iniziò così a
orbitare intorno al corpo a corpo di \emph{The Last of Us} (parte I e II).
Sul canale della GDC esistono diversi talk che raccontano cosa è stato fatto
e come, con successi, insuccessi e costi (ad esempio del \emph{motion
capture}): da quei talk si ricavarono moltissime informazioni utili.

I competitor possono essere:
\begin{itemize}
  \item \textbf{diretti}: un gioco molto simile al proprio;
  \item \textbf{di mercato}: un gioco che domina il mercato in cui ci si vuole
        affacciare;
  \item \textbf{indiretti}: un gioco con un sistema simile al proprio, ma
        tutto il resto diverso.
\end{itemize}
Regola d'oro: \textbf{là fuori, nascosto nell'ombra, c'è sicuramente un gioco
molto simile al vostro}. Va solo trovato. E se per qualche ragione non ha
venduto, bisogna capire perché.

\subsection{Lavorare in engine e fare tuning}

Il level designer lavora in engine, così come il technical designer; ma
qualsiasi game designer, a seconda di ciò che deve fare, può finire in
engine, soprattutto per il \textbf{bilanciamento} e il \textbf{tuning}.

Il \emph{tuning} è l'azione di bilanciare con precisione i dati. Esempio: in
\emph{Super Mario} si modifica quanto in alto salta il personaggio, con che
velocità, in quanto tempo raggiunge l'apice, e così via. Questo lavoro rende
il gioco:
\begin{itemize}
  \item ad alto livello, dotato della \textbf{progressione} giusta;
  \item a basso livello, \textbf{piacevole da giocare}.
\end{itemize}
Il \emph{fine tuning} è un lavoro di dettaglio, paziente e minuzioso.

\subsection{Bilanciare}

Il bilanciamento, emerso tardi nel brainstorming, è una delle attività
centrali del designer e ha anche una specializzazione dedicata
(l'\emph{economy designer}, sezione~\ref{sec:tipi-designer}).

\subsection{Analizzare il mercato}

In pre-produzione, di solito, sì: capire le tendenze, i generi che cambiano
nel tempo, la compatibilità del gioco con il momento storico, in vista ad
esempio di un possibile publisher. Va capito quanto, come e con quanto tempo
farlo, perché si è sempre limitati da molte cose.

\subsection{Sintesi}

\begin{quote}
Il designer parte dall'idea ma non si ferma lì: la confuta, lavora sui
dettagli, bilancia, fa \emph{fine tuning}; lavora con il team e fa da
\emph{manager} per diversi dipartimenti; e naviga costantemente tra i limiti.
\end{quote}

% ---------------------------------------------------------------------
\section{I limiti del design}
\label{sec:limiti}

I limiti ci sono sempre, e nel mondo del design sono spesso numerosi. I più
importanti entro cui navigare sono:
\begin{description}
  \item[Tempo] ad alto livello, l'idea deve essere realizzabile nel tempo di
        produzione disponibile; a basso livello, il proprio lavoro
        settimanale deve essere fattibile nel tempo che si ha. Non ci si può
        dividere tra mille feature team se realisticamente se ne può seguire
        uno solo.
  \item[Tecnologia] si può o non si può fare qualcosa? Ci sono i server per un
        multiplayer con un certo numero di giocatori? Ci si può permettere,
        ad esempio, una fauna che reagisce alla fisica come quella di
        \emph{GTA VI}, che richiede il meglio dei \emph{technical artist}?
  \item[Background e competenze] che cosa sa fare il team? Realisticamente,
        chi ha lavorato dieci anni su giochi horror viene assunto per un
        horror, perché ha le competenze per risolvere problemi già visti.
  \item[Persone e motivazione] il team ha la motivazione e le competenze per
        fare un certo tipo di gioco, e per lavorare insieme?
  \item[Denaro] tutto ruota intorno ai soldi. Che si lavori da soli o in mille,
        le persone vanno pagate e il gioco deve avere un ritorno economico.
        Bisogna decidere dove tracciare la linea: raggiungere il
        \textbf{break-even} (guadagnare esattamente quanto speso) o puntare a
        grandi guadagni. Con un \textbf{publisher}, inoltre, una parte
        consistente dei ricavi va a lui (ripartizioni dell'ordine di 80/20 o
        70/30, che dipendono dal contratto firmato).
\end{description}

\subsection{Non siete ``gli uomini delle idee''}

Per qualcuno può essere un problema scoprirlo, ma il designer non è
``l'uomo (o la donna) delle idee'': accanto alla parte creativa c'è una parte
tecnica, e si è fortemente limitati da molti fattori. Questi limiti, però,
sono un aiuto:
\begin{itemize}
  \item riducono il ventaglio di possibilità potenzialmente distruttive per il
        team (cose che non si riescono a fare, che non piacciono o per cui
        mancano le risorse);
  \item aiutano a mirare bene, trovando soluzioni creative in uno spazio più
        piccolo e più attuabile.
\end{itemize}

\begin{quote}
\centering
\textbf{Le idee sono nulla, l'implementazione è tutto.}
\end{quote}

% ---------------------------------------------------------------------
\section{Ruoli e specializzazioni}
\label{sec:tipi-designer}

\subsection{La scala gerarchica}

\begin{description}
  \item[Creative Director (CD)] mantiene la \textbf{visione} e dà l'idea, in
        modo che tutto segua una certa direzione. In teoria può diventarlo
        chiunque, ma con un'esperienza notevole, nell'ordine dei decenni: si
        assume che in quel tempo abbia visto come lavorano molti dipartimenti.
        Naturalmente un CD con un background artistico ha meno competenze di
        design, e viceversa.
  \item[Director di disciplina] uno per dipartimento: \emph{game director},
        \emph{level director}, \emph{art director}, ecc.
  \item[Lead] il capo del dipartimento, che gestisce i gruppi di persone. Chi
        lavora in un dipartimento risponde al proprio lead (ad esempio un game
        designer risponde al \emph{lead designer}, non al game director).
  \item[Senior / Principal] i profili più esperti (a seconda dell'azienda).
  \item[Junior] il profilo di inizio carriera.
  \item[Intern] lo stagista.
\end{description}
Lo schema si replica uno a uno per ogni dipartimento: arte, suono, animazione
(a volte arte e animazione insieme, o \emph{rigging} e animazione insieme),
programmazione, UX/UI e così via. A fare la differenza, a parità di ruolo, è
il \textbf{livello di seniority}.

\subsection{UX e UI}

UX e UI stanno a metà tra game design e arte: a seconda dell'azienda il ruolo
si chiama \emph{UX/UI Artist} o \emph{UX/UI Designer}.
\begin{description}
  \item[UX (User Experience)] come una persona interagisce con le cose. Esempio:
        su uno smartphone, un pulsante accende e spegne lo schermo; per
        digitare la password si scorre verso l'alto e compaiono i punti da
        toccare.
  \item[UI (User Interface)] come sono ``vestite'' queste cose: ad esempio
        l'aspetto della barra dei punti vita (HP), realizzata con asset 2D o
        3D.
\end{description}

\subsection{Le specializzazioni del game design}

Più l'azienda è grande, più la specializzazione è puntuale.
\begin{description}
  \item[Gameplay / Core mechanics designer] (sono la stessa cosa) si occupa
        delle meccaniche e delle 3C.
  \item[System designer] si occupa dei sistemi e delle loro interazioni.
  \item[Economy designer] si occupa del bilanciamento; lavora soprattutto con
        fogli di calcolo (Excel, Google Sheets).
  \item[Combat designer] si occupa del combattimento; negli annunci di lavoro
        indica quasi sempre un combattimento \emph{action}, con gameplay
        reattivo (non strategico), anche se non è una regola.
  \item[Multiplayer designer] si occupa delle interazioni multigiocatore.
  \item[Technical designer] fa da ponte tra game design e programmazione.
\end{description}

\subsection{Le specializzazioni del level design}

\begin{description}
  \item[Encounter designer] posiziona i nemici nella mappa in modo che gli
        scontri siano bilanciati e la progressione equilibrata come
        desiderato.
  \item[World / map designer] popola il mondo di gioco.
  \item[Quest designer] popola il mondo con le missioni e pensa alla
        progressione al loro interno, cioè a dove il giocatore deve andare (si
        vedano ad esempio gli annunci di lavoro di CD Projekt Red o di
        Larian).
  \item[Technical level designer] fa da ponte tra level design e
        programmazione.
\end{description}

\subsection{Discipline vicine: suono, UI, produzione}

Il monte ore del corso non permette di uscire come \emph{narrative designer}
(o \emph{sound designer}) completi. Ciò non toglie che si possano coltivare
queste discipline:
\begin{itemize}
  \item nei progetti di gruppo si può occuparsi anche del suono, della UI/UX,
        o provare ad applicare metodologie come Agile per fare
        \emph{producing} e gestire il team;
  \item il tempo dedicato a queste attività va però sottratto al game design,
        e dipende da quante persone ci sono nel team, da quanto lavoro c'è e
        da quanto è grande il progetto;
  \item avere sviluppato e rilasciato un progetto, anche piccolo (anche una
        \emph{game jam}), documentando i propri processi, è già un ottimo
        argomento da presentare in futuro;
  \item la via migliore resta lavorare nell'industria e, da lì, chiedere di
        imparare anche altre discipline.
\end{itemize}
Un gioco si può fare su qualsiasi cosa: sugli scarafaggi, sui piedini delle
sedie, sui pennarelli. Si parte avvantaggiati da ciò che è vicino alla propria
sensibilità: chi ama la musica può pensare a un gioco con meccaniche legate
alla musica, con un gameplay innovativo ma contenuto e realizzabile. Il design
è un cappello che comprende molte discipline.

% ---------------------------------------------------------------------
\section{Approfondimenti emersi in lezione}
\label{sec:approfondimenti}

\subsection{Roguelike e roguelite}

Un esempio di come il marketing conia generi: Supergiant Games, con
\emph{Hades}, rese celebre il termine \textbf{roguelite} al posto di
\emph{roguelike} (pur non essendo il primo roguelite in assoluto). La
differenza:
\begin{itemize}
  \item il \textbf{roguelike} si fonda su una serie di punti fissati in una
        sorta di \emph{manifesto} che definisce che cosa rende un gioco
        roguelike; uno su tutti: a ogni morte si ricomincia la \emph{run} da
        capo;
  \item il \textbf{roguelite} rompe alcuni di questi punti, in particolare
        introducendo una \textbf{progressione esterna alla run}: il giocatore
        può potenziarsi fuori dalla singola partita per vincere più
        facilmente in quelle successive. \emph{Hades} ne è il caso principe.
\end{itemize}
Analogamente, un mobile game sviluppato dal direttore del corso fu
promosso come ``il primo \emph{donation game} al mondo'': un gioco che
convertiva il tempo di gioco in donazioni reali a enti benefici per animali.
Anche qui un neologismo creato lato marketing per definire un genere.

\subsection{Applicazioni fuori dall'industria videoludica}

Le competenze sviluppate non servono solo ai videogiochi:
\begin{itemize}
  \item i \textbf{render architettonici} oggi possono essere realizzati con
        Unreal, e diversi professionisti formati nel game design lavorano in
        questo campo;
  \item la \textbf{gamification}, cioè l'applicazione di regole videoludiche
        a contesti non ludici (mostre interattive, progetti culturali,
        consulenze come quella per applicare questi concetti alle macchine da
        allenamento di un'azienda di fitness), è in Italia uno dei modi più
        redditizi di spendere queste competenze, spesso con finanziamenti
        pubblici.
\end{itemize}

% ---------------------------------------------------------------------
\section{Il primo task}
\label{sec:primo-task}

\begin{enumerate}
  \item Creare un indirizzo e-mail \textbf{professionale} (consigliato: un indirizzo Gmail nella forma
        \texttt{nome.cognome@gmail.com}; l'anno di nascita è accettabile).
        Va bene anche un indirizzo professionale già esistente usato per
        lavoro. Da evitare indirizzi poco seri (``non \texttt{blackfire67}'').
  \item Una persona della classe raccoglie tutti gli indirizzi in un unico
        documento e lo invia via e-mail al docente, che aggiungerà tutti al
        Drive condiviso del corso (dove verranno caricate lezioni, slide e
        materiali, per avere i documenti tutti nello stesso posto).
  \item Installare Discord e iscriversi al server del corso, dove passano le
        comunicazioni ufficiali. È consuetudine, e utile, creare anche un
        gruppo tra studenti.
\end{enumerate}
\textbf{Scadenza}: prima della lezione del venerdì.

% ---------------------------------------------------------------------
\section*{Punti chiave del capitolo}
\addcontentsline{toc}{section}{Punti chiave del capitolo}

\begin{itemize}
  \item Il corso è una \textbf{simulazione di lavoro}: ci si tratta da
        colleghi, con una gerarchia da studio di sviluppo.
  \item Si impara con l'\textbf{inquiry-based learning}: lezione, task,
        dubbi, domande, condivisione, progetti.
  \item Il fallimento affrontato con la mentalità giusta è la principale fonte
        di crescita; i feedback riguardano il lavoro, mai la persona.
  \item Il lavoro del designer è per metà creativo e per metà
        \textbf{ricerca e metodo scientifico}; il game design non è una
        scienza esatta.
  \item \textbf{Documentazione} e \textbf{pitching} sono il cuore del
        mestiere.
  \item Il designer progetta meccaniche e sistemi (partendo dalle
        \textbf{3C}), documenta, media, prototipa, confuta le idee, studia i
        competitor, bilancia e fa tuning.
  \item La leadership del designer è \textbf{al servizio} del team, non
        verticale.
  \item I limiti (tempo, tecnologia, competenze, persone, denaro) restringono
        e indirizzano la creatività.
  \item \textbf{Le idee sono nulla, l'implementazione è tutto.}
\end{itemize}
